\section{Discussion}
\label{sec:discussion}

\textbf{What the evidence supports.} The negative result is robust: across the period sweep, the factorial and three tabular backbones, a schedule fixed before training synchronises the wrong tier at the wrong time, and the bound says why. The positive result is conditional on the tables: at a given quota, \tierfed{} is or is not non-inferior to FedPer on AUROC and on AUPRC according to the $p$-values generated from the logs, and the abstract and Section~\ref{sec:results-fixed} state whichever is the case. We ask the reader to weigh three columns in particular: the AUPRC non-inferiority test, because AUPRC is the appropriate metric at this prevalence and an AUROC-only claim would be incomplete; the ``differentiated'' column of Table~\ref{tab:budget}, because it says whether the divergence-driven mechanism is doing anything beyond a uniform period at that quota; and the window-compliance column of Table~\ref{tab:network}, because it is the quantity an operator will actually check.

\textbf{Where \tierfed{} does not help.} Any backbone or quota at which Table~\ref{tab:arch} or Table~\ref{tab:budget} shows a deficit against FedPer is reported as such; we do not restrict the claim to the settings that work. On a deep tabular ResNet the three-tier partition may be too coarse for the way residual connections redistribute divergence, and a per-block partition is the natural remedy. Under DP, parameter frugality and noise concentration pull in opposite directions and the table decides which wins. Where the fixed schedule or FedLAMA-style intervals reach a higher \emph{peak} AUROC than their final value, the early-stopped column shows it, and an operator who can afford validation-based stopping should read that column.

\textbf{Limitations.}
(i) The sites of the tabular benchmark are specialty buckets with a share cap, not hospitals; the eICU table with true hospital identifiers is the proper test and is rendered only when the credentialed data were present at build time.
(ii) The bound is for SGD and the benchmark uses Adam; the SGD arm in the supplement checks the ordering but does not make the bound apply to Adam.
(iii) $\Gh_g$ is a dispersion proxy that scales with the interval that produced it; the controller is a feedback loop on that proxy and we have not proved its stability, only measured its behaviour.
(iv) The network model omits scheduling, fading and retransmission, and its profiles are representative targets rather than measurements.
(v) The inversion attack is generic and the DP budgets are large; neither result should be read as a privacy guarantee.
(vi) Full participation is assumed; partial participation adds the usual sampling term to (\ref{eq:bound}) and would permit subsampling amplification in (\ref{eq:rdp}).
(vii) Ten seeds on two public datasets and one imaging collection is a benchmark, not a clinical validation.

\textbf{Practical guidance.} For a tabular backbone of this size, an operator should set $\rho$ from the window and slice budget through (\ref{eq:compose}) and (\ref{eq:netcost}), keep $T_w=5$ and $K_{\max}=8$, compose with 8-bit quantisation, and read Table~\ref{tab:budget} for the accuracy cost of the chosen quota on this benchmark before transferring the setting to another one.
